All displayed masses in the latent allocation are nonnegative and sum to
\[
0.075+0.10+0.075+0.25+0.50=1.
\]
Every unit has \((D_0,D_1)=(0,1)\), so every unit is a complier and there are no defiers.  The only selection types are \((S_0,S_1)=(1,1),(0,1),(0,0)\), so weak selection monotonicity holds with \(d(x^\star)=+1\).  Taking the fair instrument independently and imposing consistency and exclusion produces, under \(Z=0\), the survivor masses \((0.075,0.10,0.075)\) at \(D=0\) and total unselected mass \(0.75\).  Under \(Z=1\), the selected outcome masses are
\[
(0.075,0.10,0.075)+(0.025,0.15,0.075)=(0.10,0.25,0.15)
\]
at \(D=1\), with unselected mass \(0.50\).  Thus the latent law belongs to \(\mathcal M\) and induces exactly \(P_{\mathrm{obs}}^\star\).

Because all units are compliers, the observable capacity contrasts are
\[
\ell=(0.075,0.10,0.075),\qquad h=(0.10,0.25,0.15).
\]
Consequently,
\[
q_0=m=0.25,\qquad q_1=0.50,\qquad \Delta q=0.25.
\]
The increasing-branch lower prefix differences are \(-0.025,-0.175,-0.25\), so the zero cut gives \(B_L=0\).  For the upper endpoint, the three threshold values \(\ell_{<t}+h_{>t}\), for \(t=0,1,2\), are
\[
0.40,\qquad 0.225,\qquad 0.175.
\]
After also taking the mass cap \(m=0.25\), this gives \(B_U=0.175\).  Hence the sharp interval is
\[
\Theta_I(P_{\mathrm{obs}}^\star)
=\left[\frac0{0.25},\frac{0.175}{0.25}\right]=[0,0.7].
\]
The displayed diagonal survivor allocation itself attains the lower endpoint.  For a transparent upper attaining survivor coupling, send masses \(0.075\) from \(Y_0=0\) to \(Y_1=1\), \(0.10\) from \(Y_0=1\) to \(Y_1=2\), and \(0.075\) from \(Y_0=2\) to \(Y_1=0\).  Its column usage \((0.075,0.075,0.10)\) is bounded by \(h\), and its benefit mass is \(0.175\); allocating the residual \(h-(0.075,0.075,0.10)=(0.025,0.175,0.05)\) to \((S_0,S_1)=(0,1)\) completes the upper law.  This also directly verifies the endpoint-attainment conclusion in this witness.